PaperSwarm's evidence layer consists of two append-only JSONL ledgers. The first, \texttt{artifacts.jsonl}, records each experiment artifact as a tuple with an artifact identifier, a path, a SHA-256 digest, a byte length, and the producer. The second, \texttt{claims.jsonl}, records each numeric claim as a tuple with a claim identifier, the artifact identifier from which the value is drawn, a JSON Pointer locating the value within that artifact, the registered value, and a display string. A claim may be registered only when the stated value equals the value actually present at the referenced JSON Pointer; any mismatch aborts the run before the entry is written. The ledgers therefore accumulate verified observations rather than asserted results.

During writing, section authors do not emit digits. They insert a result macro that takes a claim identifier, and the assembler later resolves that identifier by substituting the verified value and display string from the claims ledger at build time. This arrangement keeps the writer separated from the raw measurements and makes every reported figure traceable to a specific location in an immutable artifact.

A gate predicate is applied to each draft text. It rejects a section when the text references an unknown claim identifier, when a referenced claim fails re-verification against its artifact, when the text redefines the result macro, or when a literal number appears in a result context. Re-verification occurs before assembly so that an artifact cannot change after registration without invalidating the dependent claim. A rejected draft is returned to the writer with an explicit list of violations and is rewritten from scratch, up to a fixed attempt budget. If the budget is exhausted, the run terminates without producing a PDF rather than shipping a weaker paper.

A separate structural check rejects a section that writes another section's material, preserving scope discipline across the multi-agent system. In the present run, the experiment is a controlled synthetic regression study, and the baseline follows the ridge-regression method due to \citet{hoerl1970ridge}. The present section reports no experimental numbers.
